\section{Plain versus robust}\label{sec:plainrobust}

The plain conjecture rules out codes that are correct on \emph{all} inputs; the robust one rules out
codes correct on a $2^{-O(nw)}$ fraction of the inputs. Robust implies plain trivially. The plain
form already refutes oblivious linear-time transposers (\cref{thm:oblivious}); if plain implied
robust, it would refute~(T). This section isolates the exact gap and proves the implication for
several natural classes of codes and sets. The general case is open.

\subsection{The gap is an oracle hint}

A \emph{hint code} with $2^s$ hint values consists of a function $h$ of the whole input and a family
of causal codes $(C_v)_v$ such that on every input $x$ the code $C_{h(x)}$ solves complete exchange
for $x$. It models an oracle that sees the whole input and tells every node $s$ bits; the hint is
uncharged (no edge carries it and no node computes it). Write $\mathrm{NoCode}^{\mathrm{hint}}(s)$
for ``no hint code with $2^{s(n)}$ values exists'' (\lean{NoCode_TM_hintBits}), with the same
quantifiers as the conjecture.

\begin{theorem}[Robust is plain with an oracle hint]\label{thm:hint}
\formal{Tape/Statements/Coding.lean}{robustNoCode_iff_hint, not_linearTimeTransposable_of_noCode_hint, hintCodes_of_linearTimeTransposable, noCode_of_noCode_hint, solves_shiftCode, exists_shift_cover, hint_of_code, code_of_hint, robustNoCode_iff_coent, noCode_iff_coent_zero, coent_mono, hintBits_of_coent, coent_of_hintBits, noCode_bits_iff}
\begin{enumerate}[label=(\alph*)]
\item On a fixed diagram and at fixed capacities: a hint code with $2^s$ values gives a code on a set
of co-entropy $s$; a code on a set of co-entropy $m$ gives a hint code with $2^m(L+1)$ values.
\item The Tape Coding Conjecture is equivalent to: for every $C$, no hint code with an $Cnw$-bit
hint exists. With a constant hint this gives the plain form. (T) yields, for every large $n$, a single
valid diagram with codes indexed by an $O(nw)$-bit hint solving complete exchange on every input.
\item The co-entropy ladder: $\mathrm{NoCode}(0)$ is the plain form, $\mathrm{NoCode}(Cnw)$ for every
$C$ is the robust form, the rungs are monotone, and hints and co-entropy agree up to $3w+1$ bits.
\item A side channel whose contents the code computes itself is free: for $s(k,n)\le C(k)\,nw$, the
plain form with $s$ extra bits on every memory edge is equivalent to the plain form.
\end{enumerate}
\end{theorem}

\begin{outline}
(a) In one direction take the most popular hint value. In the other, a code for $F$ is a code for
every shift $F\oplus z$ (each node XORs its visible rows with $z$, and each output bit is XORed back),
and $2^m(L+1)$ shifts of $F$ cover every input, by counting with Bernoulli's inequality. (b) and (c)
follow from (a). (d) Extra memory bits become stationary tapes, as in \cref{thm:profile}.
\end{outline}

So the entire gap between plain and robust is whether an $O(nw)$-bit oracle broadcast can be
replaced by one that the network computes itself. For the application to~(T), co-entropy $O(n)$
suffices (\cref{thm:coent-linear}), so only hints of $O(n)$ bits matter; this is still $\omega(1)$
bits. The first rung of the ladder is already open: does ``no code on all inputs'' imply ``no code
on half of the inputs''? This is a one-shot question of zero error versus constant error for this
class of networks.

\subsection{Partial results}

Throughout this subsection we assume the plain conjecture \lean{NoCode_TM} and conclude, for all
constants and all large $n$, that certain codes do not exist at the profile capacities. The
constructions prepend or append computations to the diagram (an input sweep, a selection phase,
copies of the work walks), so they change $k$ and $c$ by constant factors; this is harmless because
the plain conjecture quantifies over all constants.

\begin{theorem}[Affine sets and additive hints]\label{thm:pr-affine}
\formal{Tape/Statements/PlainRobust.lean}{PlainRobust.noCode_affine_of_noCode, PlainRobust.noCode_of_affine_subset, PlainRobust.noCode_additive_of_noCode, PlainRobust.noCode_additiveCover_of_noCode, PlainRobust.noCode_affine_exc_of_noCode}
Under the plain conjecture, no code exists on a set $F$ that contains a nonempty affine set
$\{x:\sum_ig_i(\text{row }i\text{ of }x)=b\}$, where the $g_i$ are $\F_2$-linear maps to $\F_2^s$ with
$s\le Cnw$; more generally, on a set $F$ such that, for an additive hint $\sum_ig_i(\text{row }i)$ of
$s\le Cnw$ bits and $K=O(1)$ shift tables, every input is moved into $F$ by one of the shifts; and on
such an affine set from which a set $E$ with $2^s\cdot2\abs{E}^2<2^L$ is removed.
\end{theorem}

\begin{outline}
Prepend an input sweep: the input head goes from $0$ to $L-1$ and back while the work heads are
parked, so every row is read in the prefix and the original diagram follows unchanged. The prefix
computes the additive hint along its spine, and the original part runs the code shifted by the shift
table entry of the hint, which it reads from the low $s$ bits of its memory edges. This is the
\emph{syndrome trick} of external-memory coding in the language of tape diagrams. The version with
$K$ shifts adds a selection among the $K$ shifted codes (\cref{thm:pr-select}), and the version with
exceptions adds voting (\cref{thm:pr-vote}).
\end{outline}

\begin{theorem}[Codes failing on few inputs]\label{thm:pr-vote}
\formal{Tape/Statements/PlainRobust.lean}{PlainRobust.noCode_cosmall_of_noCode}
Under the plain conjecture, no code is correct on all inputs outside a set $E$ with
$2\abs{E}^2<2^L$.
\end{theorem}

\begin{outline}
Triple every work walk; the block network is unchanged and the profile capacities triple. Run three
shifted copies of the code, with shifts $0,z_1,z_2$ such that $z_i\oplus z_j\notin E\oplus E$, and
output the majority of the three decoded bits: on every input at most one copy errs.
\end{outline}

\begin{theorem}[Selection among candidates]\label{thm:pr-select}
\formal{Tape/Statements/PlainRobust.lean}{PlainRobust.noCode_hintConst_of_noCode, PlainRobust.noCode_iff_hintConst, PlainRobust.noCode_netCover_of_noCode, PlainRobust.noCode_netHint_of_noCode}
Under the plain conjecture, no hint code with a constant number of hint bits exists; equivalently, the
plain conjecture is unchanged by a constant oracle hint. More generally, no code exists on a set $F$
for which a \emph{network-computable} hint (computed by a causal code within the profile capacities of
a linear-length prefix diagram) together with $O(1)$ shift tables moves every input into $F$; and no
family of codes indexed by a network-computable hint of $O(nw)$ bits solves complete exchange on every
input.
\end{theorem}

\begin{outline}
Run the $K$ candidate codes side by side (clamped and packed). Each node compares the candidates'
would-be output bits at its own output positions and, for every pair of candidates without a recorded
difference, records one position where they differ, with the first candidate's bit; this state
travels along the spine. An appended input sweep fetches the true bit at every recorded position, and
a fresh copy of the diagram runs a candidate that never lost a comparison. A correct candidate never
loses, and every incorrect one loses against it, so the winner is correct. This is discrimination,
not verification: telling two candidates apart takes one fetched bit. For network-computable hints,
concatenate the prefix computation with the diagram and carry the hint along the spine.
\end{outline}

\begin{theorem}[Linear and nearly linear codes]\label{thm:pr-linear}
\formal{Tape/Statements/PlainRobust.lean}{PlainRobust.noCode_linear_of_noCode, PlainRobust.noCode_linearExcept_of_noCode, PlainRobust.noCode_linearPlus_of_noCode, PlainRobust.noCode_linearPlusImage_of_noCode}
Under the plain conjecture, no $\F_2$-linear code solves complete exchange on a set of co-entropy at
most $Cnw$. The same holds for codes that are linear outside a set of edges carrying $O(nw)$ bits, for
codes that are linear plus nonlinear corrections of $O(nw)$ bits in total, and for codes that are
linear plus corrections taking at most $2^{O(nw)}$ joint values.
\end{theorem}

\begin{outline}
For a linear code the set of inputs on which it is correct is closed under XOR, hence a subspace; if it
contains a set of co-entropy $Cnw$ it has codimension at most $Cnw$, so it is the kernel of an additive
syndrome map, and \cref{thm:pr-affine} applies. For the nearly linear codes, the nonlinear values form
an internally computed hint of $O(nw)$ bits: replacing them by constants gives an affine code, a
$2^{-O(nw)}$ fraction of the good inputs share their nonlinear values by pigeonhole, and XOR-ing these
inputs with one of them gives inputs on which the linear part is correct.
\end{outline}

\begin{theorem}[Bounded-degree constraints and decompositions]\label{thm:pr-degree}
\formal{Tape/Statements/PlainRobust.lean}{PlainRobust.noCode_lowDegree_of_noCode, PlainRobust.noCode_quadForm_of_noCode, PlainRobust.noCode_decomp_of_noCode, PlainRobust.noCode_crossQuad_of_noCode}
Under the plain conjecture, no code exists on a set containing the zero set of a Boolean function of
degree at most $D$ (for constant $D$, and the zero set nonempty), in particular the zero set of a
quadratic form. No code exists on a set that, for a partition of the input positions into $K=O(1)$
parts, contains every input supported on a single part; for instance the common zero set of any
family of quadratic forms each of whose monomials joins two different parts.
\end{theorem}

\begin{outline}
The zero set of a function of degree at most $D$ is covered by $2^D$ translates of itself: if a
derivative of order $D'\le D$ is nonzero at $x$, some translate of $x$ by a sum of a sub-list of the
directions lies in the zero set, by a parity argument; so selection among $2^D$ shifted codes applies.
For a decomposition, run $K$ copies of the code, copy $j$ on the input masked outside part $j$, which
lies in $F$, and XOR the outputs: transposition is linear.
\end{outline}

\begin{theorem}[Errors]\label{thm:pr-errors}
\formal{Tape/Statements/PlainRobust.lean}{PlainRobust.noCode_fewErrors_of_noCode, PlainRobust.noCode_sparseErrors_of_noCode, PlainRobust.noCode_sparseEntryErrors_of_noCode, PlainRobust.noCode_fewErrorsCover_of_noCode}
Under the plain conjecture, every code that stays within the profile capacities on every input has
more than $2^{Hnw}$ distinct error vectors $C(x)^{\T}\oplus x$, for every constant $H$. In particular,
no such code gets all but at most $An$ output bits right on every input (transposition with $O(n)$
wrong bits is as hard as exact transposition); the same holds for $O(n)$ wrong entries, and for a
choice among $O(1)$ such codes per input.
\end{theorem}

\begin{outline}
Syndrome correction. Fix a check matrix $\Phi$ with $t=2Hnw+1$ rows that is injective on the possible
error vectors (a random one is). Run the code; every output bit equal to $1$ adds the corresponding
column of $\Phi$ on the spine; an appended input sweep adds the column of every input bit equal to $1$,
so the spine carries the syndrome $\Phi e(x)$; a second run of the code XORs each output bit with the
error bit decoded from the syndrome. The $t$ extra bits are paid for by stationary tapes.
\end{outline}

\begin{theorem}[Column covers]\label{thm:pr-columns}
\formal{Tape/Statements/PlainRobust.lean}{PlainRobust.noCode_colCover_of_noCode, PlainRobust.noCode_rowMapCover_of_noCode, PlainRobust.noCode_colAggregate_of_noCode}
Under the plain conjecture, no code exists on a set $F$ for which $K=O(1)$ locally computable maps
$\Phi_\kappa$ into $F$ exist, each keeping the input columns of its class $\kappa$ (the classes
partitioning the columns); in particular for maps acting on rows. For even $n$, no code exists on a set
containing $\{x:\sum_jv(\text{column }j\text{ of }x)=0\}$ for a fixed $v$ with values in $\F_2^m$.
\end{theorem}

\begin{outline}
Run $K$ copies, copy $\kappa$ on $\Phi_\kappa(x)$ (each node computes the copy's visible input from its
own), and read column $j$ from the copy of its class. For the aggregate constraint, the inputs
$[x_{\mathrm{left}}\mid x_{\mathrm{left}}]$ and $[x_{\mathrm{right}}\mid x_{\mathrm{right}}]$ satisfy it,
since each term appears twice.
\end{outline}

\subsection{Why the general case is hard}

We record, without formal proofs, why these methods stop short of the general implication.
\begin{itemize}[leftmargin=2em]
\item \emph{A general hint-removal theorem must use the code, not only the size of $F$.} A counting
argument over (automaton, shift table) pairs shows that a random set of density $2^{-m}$ has no
streaming covering hint of $O(nw)$ bits.
\item \emph{A hint is cheap in information but not in flow.} An $s$-bit broadcast tells at most $s$
bits about all outputs together, but delivering it to $\Theta(n)$ nodes costs $\Theta(L)$ bits of
capacity in aggregate; edge-removal style arguments do not apply as they stand. The first rung of the
ladder is a one-shot instance of zero error versus constant error, related to the edge-removal problem
and to the continuity of network capacity under dependent sources~\cite{LangbergEffros2021}.
\item \emph{The machine-side view.} Plain would refute (T) if (T) implied oblivious (T), or merely a
linear-time transposer whose block-level trace does not depend on the input. Codes of trace classes are
self-detecting: at every slot the simulation knows whether the run still follows the popular trace. A
hint-removal argument restricted to self-detecting codes would suffice; none is known.
\end{itemize}
